\documentclass{article}
\usepackage{graphicx}
\usepackage{amsmath}
\usepackage{amssymb}
\usepackage{booktabs}
\usepackage{hyperref}
\usepackage{geometry}
\geometry{margin=1in}
\usepackage{float}
\usepackage{subcaption}

\title{DIME - Distributed Infrastructure Management Environment}
\author{Naseer Hussain \and Shivangi Sharma \and Nithish Sri Ram}
\date{April 2026}

\begin{document}

\maketitle

\begin{abstract}
We introduce DIME (Distributed Infrastructure Management Environment), a simulated
8-node Kubernetes cluster in which a language model acts as an autonomous
site-reliability engineer. The agent observes per-node CPU, memory, queue depth, and
tail-latency telemetry at each timestep and must issue a single management command.
We apply Group Relative Policy Optimization (GRPO) to Qwen3-8B with a four-component
composite reward combining structural compliance, syntactic validity, physics-based
environment feedback, and a deterministic triage oracle. We document several failure
modes encountered during development --- including a reward cliff that caused
zero-variance advantages and complete training collapse, a priority inversion in the
reference policy that caused the model to learn a plausible but suboptimal behavior,
and a generation length truncation issue that silently suppressed all learning signal ---
and show how each was diagnosed and corrected. After 44 minutes of fine-tuning on a
single A100-SXM4-80GB, the model improves from a zero-shot average of 0.3946 to 0.4649
across 14 failure-scenario tasks (+17.8 pp, +44.8\% relative). We release the
environment, training code, and fine-tuned model checkpoint to support future work on
RL-trained infrastructure agents.
\end{abstract}

\section{Introduction}

Autonomous management of distributed infrastructure is among the most operationally
demanding tasks an AI agent can face. Unlike text generation or code synthesis, the
SRE setting imposes hard real-time constraints: structured, parseable output on every
step; action budgets that deplete permanently; cascading failure dynamics that amplify
early mistakes; and partial observability when node telemetry drops out. Errors do not
produce a wrong answer in a response --- they take down a production database.

Recent work has demonstrated that LLMs can act as useful tools in cloud operations
contexts. RCAgent [Wang et al., 2024] augments LLMs with tool-calling for root cause
analysis in Alibaba Cloud's production systems. AIOpsLab [Chen et al., 2025], a
framework from Microsoft Research presented at MLSys'25, provides a holistic evaluation
harness for AI agents in cloud operations across fault injection, metric analysis, and
remediation tasks. KubeIntellect [Seyedkazemi Ardebili \& Bartolini, 2025] achieves
93\% tool-synthesis success on natural-language Kubernetes queries through a modular
LLM orchestration architecture. However, these systems either rely on prompting or tool
augmentation without policy optimization, or evaluate on single-turn query answering
rather than multi-step episode control. Fine-tuning an LLM's weights to act as a
closed-loop controller in a simulated infrastructure environment --- and documenting
what breaks along the way --- is the focus of this paper.

The broader landscape of LLM agent benchmarks (WebArena [Zhou et al., 2023], SWE-bench
[Jimenez et al., 2024], ALFWorld [Shridhar et al., 2021]) has established that zero-shot
performance of capable frontier models is a reasonable baseline; we adopt the same
before/after evaluation methodology. What sets the infrastructure domain apart from web
navigation or code editing is the \textit{continuous control} nature of the task: the
agent must maintain system health across dozens of sequential decisions, and a single
mis-prioritized action (throttle instead of restart) can lock an episode into a
suboptimal attractor state.

We make the following contributions:

\begin{enumerate}
    \item \textbf{DIME}: a configurable, reproducible 14-task Kubernetes simulation
    benchmark with partial observability, error budget constraints, and
    Prometheus-style telemetry.

    \item \textbf{Reward redesign}: we replace an early hard-cliff penalty (uniform
    large negative reward on DB failure) with a seven-component differentiable signal bounded to $[-5, +5]$, and
    show that the original formulation would have made GRPO (or any policy gradient
    method) impossible to apply regardless of hyperparameter tuning.

    \item \textbf{Oracle priority analysis}: we identify and correct a priority
    inversion in the reference triage policy that caused the model to learn
    \texttt{throttle} in database-failure states rather than \texttt{restart\_node(0)},
    and quantify the downstream impact (+0.044 benchmark score).

    \item \textbf{Failure documentation}: we describe five training configurations that
    failed before convergence, each diagnosable from standard TRL metrics, providing a
    practical troubleshooting guide for GRPO on environment-coupled reward functions.

    \item A \textbf{public model checkpoint} at
    \texttt{Naseer-010/Qwen3-8B-Finetuned-DIME} on the Hugging Face Hub.
\end{enumerate}

\section{Related Work}

\textbf{LLM agents for cloud and infrastructure operations.} RCAgent [Wang et al.,
2024] demonstrates LLM-driven root cause analysis in a production Alibaba Cloud
environment, using tool-calling rather than weight updates. AIOpsLab [Chen et al.,
2025] is the closest antecedent to DIME: it provides a systematic evaluation framework
for AI agents in cloud incident management, covering detection, diagnosis, and
mitigation. Our work differs in that we \textit{fine-tune} the model's weights via
GRPO rather than prompting, and focus on the episodic closed-loop control problem.
Concurrent work on autonomous resource management in microservice systems [Zou et al.,
2025] uses tabular RL (PPO) on environment simulators; we target an LLM policy directly.

\textbf{RL fine-tuning of language models.} GRPO [Shao et al., 2024], introduced in
the DeepSeekMath paper, eliminates the need for a separate value network by using
group-internal advantage normalization. PPO [Schulman et al., 2017] is the dominant
prior approach but requires a critic network that doubles memory usage and training time
at LLM scale. We choose GRPO specifically because the CPU-bound reward evaluation (each
step requires a full environment simulation) makes small batch sizes necessary, and the
absence of a critic is a meaningful practical advantage.

\textbf{ReAct and reasoning-action integration.} ReAct [Yao et al., 2023] established
the pattern of interleaving natural language reasoning with action emission that we adopt
in our output format. Our \texttt{<reasoning>...</reasoning><action>...</action>}
scaffold follows this paradigm and is reinforced by the format reward component.

\textbf{Parameter-efficient fine-tuning.} We use LoRA [Hu et al., 2022] to limit the
trainable parameter count to 1.05\% of total weights (349 MB vs.\ 16 GB base model),
making the fine-tuned checkpoint cheap to distribute and easy to swap over the base
model.

\section{The DIME Environment}

\subsection{Cluster Model}

DIME simulates an 8-node cluster. Node 0 is a stateful PostgreSQL-style database and
acts as a single point of failure (SPOF). Nodes 1--7 are stateless application workers.
At each step $t$ the agent receives observation

\[
s_t = \bigl(\mathbf{c} \in [-1,1]^8,\; \mathbf{m} \in [0,1]^8,\; \mathbf{q} \in \mathbb{Z}_{\geq 0}^8,\; F \subseteq \{0,\ldots,7\},\; \lambda,\; \lambda_{99},\; w_{\mathrm{io}},\; B,\; t\bigr)
\]

where $\mathbf{c}$ are per-node CPU utilizations ($c_i = -1$ denotes telemetry dropout),
$\mathbf{m}$ are memory utilizations, $\mathbf{q}$ are pending request queue depths,
$F$ is the set of failed node indices, $\lambda$ is mean latency (ms), $\lambda_{99}$
is P99 tail latency, $w_{\mathrm{io}} \in [0,1]$ is a DB disk I/O saturation proxy,
and $B \in \mathbb{Z}_{\geq 0}$ is the remaining error budget for load-shedding actions.

\subsection{Action Space}

The agent emits one action per step from six types:

\begin{table}[h]
\centering
\begin{tabular}{lll}
\toprule
\textbf{Action} & \textbf{Effect} & \textbf{Constraint} \\
\midrule
\texttt{restart\_node($i$)} & Restores failed node $i$ & 2-step boot delay; 5-step cooldown \\
\texttt{reroute($i \to j$)} & Shifts load fraction from $i$ to $j$ & Both nodes must be alive \\
\texttt{scale\_up} & Adds a capacity node for 10 steps & Costs one budget unit \\
\texttt{throttle($r$)} & Sets ingress acceptance rate $r \in [0,1]$ & Burns error budget \\
\texttt{query\_logs($i$)} & Resolves telemetry dropout on node $i$ & --- \\
\texttt{no\_op} & Passive observation & --- \\
\bottomrule
\end{tabular}
\caption{Action space of the DIME environment.}
\end{table}

Actions are expressed as natural-language kubectl strings (e.g.,
\texttt{kubectl delete pod node-3}), which DIME parses into the structured action types
above. Malformed commands receive a \texttt{ParseError} penalty but do not crash the
episode.

\begin{figure}[H]
  \centering
  \includegraphics[width=\textwidth]{fig9_architecture.png}
  \caption{%
    DIME cluster architecture and agent loop. Node-0 (red, PostgreSQL DB) is the
    single point of failure; nodes 1--7 are stateless workers. At each step the LLM
    agent receives observation $s_t$ from the environment and emits action $a_t$ as a
    natural-language \texttt{kubectl} string. The composite reward $R(a_t, s_t)$ is
    computed internally and used only during training.
  }
  \label{fig:architecture}
\end{figure}

\subsection{Request Simulation}

Requests are simulated as a blend of 60\% real production trace data (Alibaba cluster
trace v2018) and 40\% stochastic simulation, mixed per step. This blend gives DIME
realistic burst patterns while maintaining episode-to-episode variability. The
stochastic component means per-task scores vary across runs even with a fixed model; we
discuss the implications for evaluation in Section~8.3.

\subsection{Benchmark Tasks}

DIME defines 14 canonical failure scenarios:
\texttt{traffic\_spike}, \texttt{node\_failure}, \texttt{cascading\_failure},
\texttt{flash\_crowd}, \texttt{level\_5\_alibaba\_trace}, \texttt{thundering\_herd},
\texttt{zombie\_node}, \texttt{memory\_leak\_slow\_burn},
\texttt{split\_brain\_io\_bottleneck}, \texttt{black\_swan\_az\_failure},
\texttt{retry\_storm}, \texttt{hot\_shard\_skew},
\texttt{connection\_pool\_deadlock}, \texttt{autoscaler\_flapping\_trap}.

Each task runs for up to 100 steps. Task score is reported as a running grader value
$\tau_t \in (0, 1]$ computed by the environment's rubric engine; the final $\tau_T$ is
the per-task score used in all tables below.

\section{Reward Engineering and Failure Modes}

\subsection{The Hard Cliff: How an Unbounded Penalty Blocked All Learning}

An early version of \texttt{ProductionSREReward} in \texttt{server/rubrics.py} returned
a hard fixed penalty whenever the database node failed:

\[
r(s, a) = \begin{cases} r_{\min} & \text{if } 0 \in F(s) \\ -c_{\mathrm{db}} \exp(\lambda_{\mathrm{db}} \cdot c_0) + \cdots & \text{otherwise} \end{cases}
\]

with $c_{\mathrm{db}} = 100$, $\lambda_{\mathrm{db}} = 4.0$.
At $c_0 \geq 0.57$ the second branch also evaluated to an extreme negative value.
Because the environment's request simulator blends 60\% real trace data with 40\%
stochastic simulation, DB CPU routinely crossed 0.57 within the first 1--3 steps of
most episodes. Once node 0 crashed, every subsequent reward was the same large
negative constant regardless of action.

This is not a hyperparameter problem --- it is a structural one. When all rewards
within a GRPO group are equal, the advantage estimate is:

\[
\hat{A}_i = \frac{r_i - \bar{r}}{\sigma_r + \epsilon} = \frac{0}{\epsilon} \approx 0
\]

No gradient flows regardless of learning rate, KL coefficient, or batch size. We
confirmed this empirically: \texttt{frac\_reward\_zero\_std = 1.0} from step 1,
\texttt{train\_loss} $\approx 0$ throughout. The same pathology would affect PPO
(value loss would fit to a constant), DPO (all preference pairs would tie), or any
policy gradient variant.

The only task that escaped was \texttt{memory\_leak\_slow\_burn} (zero-shot score
0.99): memory leaks build slowly enough that the cliff was rarely triggered, giving the
model 58 of 60 steps with informative gradients. This single outlier, where the
zero-shot model already achieved near-perfect performance, was the evidence that the
model had latent SRE capability --- the reward function was the bottleneck, not the
model.

\begin{figure}[H]
  \centering
  \includegraphics[width=\textwidth]{fig5_reward_cliff_bounded.png}
  \caption{%
    \textbf{Left:} Original reward function. At $c_0 \geq 0.57$ the signal
    collapses to a uniform large negative value, making all GRPO group advantages identically zero
    ($\hat{A}_i = 0/0$). The shaded region covers the majority of episode states
    in practice. \textbf{Right:} Redesigned reward bounded to $[-5,+5]$. Every
    point on the curve carries a non-zero gradient; the DB-failure penalty is
    $-5$ (finite and distinguishable from all other states).
  }
  \label{fig:reward-cliff}
\end{figure}

\subsection{Fixed Environment Reward $R_{\mathrm{env}}$}

The fixed \texttt{ProductionSREReward} replaces the cliff with seven differentiable components:

\[
R_{\mathrm{env}}(s, a) = \mathrm{clip}\!\left(\, r_{\mathrm{topo}} + r_{\mathrm{shed}} + r_{\mathrm{mem}} + r_{\mathrm{fric}} + r_{\mathrm{lat}} + r_{\mathrm{up}} + r_{\mathrm{eff}},\; -5,\; +5 \right)
\]

\textbf{Asymmetric topology term} (DB health):

\[
r_{\mathrm{topo}}(s) = \begin{cases}
-5.0 & \text{if } 0 \in F(s) \\
-c_{\mathrm{db}} \exp(\lambda_{\mathrm{db}} \cdot c_0) - c_w \displaystyle\sum_{i=1}^{7} c_i^2 & \text{otherwise}
\end{cases}
\]

with $c_{\mathrm{db}} = 1.5$, $\lambda_{\mathrm{db}} = 2.0$, $c_w = 0.05$.
At full DB load ($c_0 = 1.0$), $r_{\mathrm{topo}} \approx -11.1$, which is clipped by
the outer bound --- the key change from the original, where the equivalent value was
$-22026$.

\textbf{Load-shedding economics:}

\[
r_{\mathrm{shed}}(s, a) = \begin{cases}
-\psi_{\mathrm{shed}} (1 - r_{\mathrm{rate}}) \left(1 + \delta_{\mathrm{shed}} \left(\tfrac{B_{\max}}{B}\right)^2\right) & \text{if } a = \texttt{throttle}(r_{\mathrm{rate}}) \\
0 & \text{otherwise}
\end{cases}
\]

with $\psi_{\mathrm{shed}} = 0.3$, $\delta_{\mathrm{shed}} = 2.0$, $B_{\max} = 100$.
This makes throttling increasingly expensive as the error budget depletes, incentivising
the agent to reserve load-shedding for genuine emergencies.

\textbf{Memory cliff:}

\[
r_{\mathrm{mem}}(s) = \sum_{i=0}^{7} \begin{cases}
-2.0 & \text{if } m_i \geq 0.98 \\
-\min\!\left(1.5,\; b_{\mathrm{mem}} \exp\!\left(\min\!\left(10, \tfrac{1}{1-m_i}\right) - 2\right)\right) & \text{if } m_i > 0.5
\end{cases}
\]

with $b_{\mathrm{mem}} = 0.5$.

\textbf{Tail-latency penalty:}

\[
r_{\mathrm{lat}}(s) = -\min\!\left(3.0,\; 1.5 \left(\frac{\max(0,\, \lambda_{99} - 50)}{200}\right)^2\right)
\]

\textbf{Uptime bonus} (dense positive signal):

\[
r_{\mathrm{up}}(s) = \frac{8 - |F(s)|}{8}
\]

\textbf{Action efficiency} (anti-spam):

\[
r_{\mathrm{eff}}(a) = \begin{cases}
-0.15 & \text{if } a = \texttt{throttle}(\cdot) \\
-0.05 & \text{if } a \neq \texttt{no\_op} \\
0 & \text{otherwise}
\end{cases}
\]

\textbf{Temporal friction} (cold-start penalty for \texttt{scale\_up}): a running
integral that decays as queue lengths drain, bounded to $[-1, 0]$.

\subsection{Composite Reward}

The full training reward is a sum of four components:

\[
R(a, s) = \underbrace{R_{\mathrm{fmt}}(a)}_{\in[-3,+3]} + \underbrace{R_{\mathrm{val}}(a)}_{\in[-2,+2]} + \underbrace{R_{\mathrm{env}}(a, s)}_{\in[-5,+5]} + \underbrace{R_{\mathrm{tri}}(a, s)}_{\in[-0.5,+1]}
\]

\textbf{Format reward} $R_{\mathrm{fmt}} \in [-3, +3]$: rewards correct XML scaffold
(\texttt{<reasoning>...</reasoning><action>...</action>}).

\[
R_{\mathrm{fmt}}(a) = \begin{cases}
+3.0 & \text{exactly one of each tag present} \\
\sum_{\text{tag}} (+0.5 \cdot \mathbf{1}_{\text{correct}} - 1.0 \cdot \mathbf{1}_{\text{missing/dup}}) & \text{otherwise}
\end{cases}
\]

\textbf{Validity reward} $R_{\mathrm{val}} \in [-2, +2]$: rewards syntactically
parseable commands.

\[
R_{\mathrm{val}}(a) = \begin{cases}
+2.0 & \text{command parses without error} \\
+1.0 & a = \texttt{no\_op} \\
-1.0 & \text{JSON present but parse fails} \\
-2.0 & \text{no action block}
\end{cases}
\]

\textbf{Environment reward} $R_{\mathrm{env}} \in [-5, +5]$: ProductionSREReward as above.

\textbf{Triage oracle reward} $R_{\mathrm{tri}} \in [-0.5, +1]$: compares the agent's
action against a deterministic reference policy $a^*(s)$ (Section~3.4). Kept
intentionally weak so $R_{\mathrm{env}}$ (physics) dominates learning.

\[
R_{\mathrm{tri}}(a, s) = \begin{cases}
+1.0 & a = a^*(s) \\
+0.5 & \text{same action type, different parameters} \\
0.0  & a = \texttt{no\_op},\; a^*(s) \neq \texttt{no\_op},\; \text{or system healthy} \\
-0.5 & \text{otherwise}
\end{cases}
\]

Total range: $R \in [-10.5, +11]$.

\begin{figure}[H]
  \centering
  \includegraphics[width=\textwidth]{fig10_reward_deepdive.png}
  \caption{%
    \textbf{Left:} Per-component reward breakdown over a 50-step
    \texttt{connection\_pool\_deadlock} episode. $R_{\mathrm{fmt}}$ and
    $R_{\mathrm{val}}$ stabilise immediately (blue/purple); $R_{\mathrm{env}}$
    (teal) recovers after the agent throttles at step 2, driving the total reward
    (green) upward to a final score of 0.9762. \textbf{Right:} Normalised
    capability radar comparing zero-shot (gray) and fine-tuned (green) across six
    sub-metrics. DB Recovery and Multi-Node Triage show the largest gains.
  }
  \label{fig:reward-components}
\end{figure}

\subsection{Triage Oracle $a^*(s)$}

The reference policy is a priority-ordered decision tree evaluated on each state $s$:

\begin{table}[h]
\centering
\begin{tabular}{clll}
\toprule
\textbf{Priority} & \textbf{Condition} & \textbf{$a^*(s)$} & \textbf{Rationale} \\
\midrule
1 & $\exists i: m_i > 0.92$ & \texttt{restart\_node($i$)} & OOM prevention \\
2 & $0 \in F(s)$ & \texttt{restart\_node(0)} & DB is SPOF \\
3 & $w_{\mathrm{io}} > 0.80$ & \texttt{throttle(0.5)} & split-brain \\
4 & $\exists i \geq 1: c_i > 0.90,\; \bar{c} < 0.60$ & \texttt{reroute($i \to \arg\min c$)} & hot shard \\
5 & $\lambda_{99} > 100$ and $\mathrm{rr} > 150$ & \texttt{throttle(0.4)} & retry storm \\
6 & $\exists i \geq 1: c_i \in [0, 0.10)$ and $\lambda_{99} > 100$ & \texttt{reroute($i \to \cdot$)} & zombie node \\
7 & $|F(s)| \geq 2$ & \texttt{throttle(0.3)} & black swan \\
8 & $c_0 > 0.80$ & \texttt{throttle(0.7)} & protect living DB \\
9 & $\bar{c}_{\mathrm{workers}} > 0.75$ and $B > 20$ & \texttt{scale\_up} & safe scaling \\
10 & --- & \texttt{no\_op} & --- \\
\bottomrule
\end{tabular}
\caption{Triage oracle priority table.}
\end{table}

\textbf{The priority inversion.} In an earlier version of the oracle, Rule 7 (Black
Swan: $|F(s)| \geq 2$) was evaluated \textit{before} Rule 2 (DB Recovery: $0 \in
F(s)$). When both conditions held simultaneously --- multiple nodes including node 0
failed --- the oracle prescribed \texttt{throttle(0.3)} rather than
\texttt{restart\_node(0)}. Throttling when the database is dead does not restart the
database; it merely reduces load onto a cluster that cannot serve requests regardless.

A model trained under this inverted oracle learned the \texttt{throttle} response in
DB-failure states with high consistency. The behavior is \textit{locally plausible} ---
it is the right response for multi-node failures when the DB is alive --- but globally
wrong when the DB is the failed node. This is a concrete case of reward hacking via
reference policy defect: the model faithfully optimized the oracle rather than the
environment, and the oracle was wrong. Correcting the priority (putting Rule 2 before
Rule 7) was the sole change between training Run 4 (0.4206) and Run 6 (0.4649), a
+0.044 improvement.

\begin{figure}[H]
  \centering
  \includegraphics[width=0.75\textwidth]{fig11_triage_tree.png}
  \caption{%
    Triage oracle decision tree (10 rules, priority order top-to-bottom).
    The \textcolor{red}{red zone} marks the original priority inversion: Rule 7
    (Black Swan) was evaluated before Rule 2 (DB Recovery), causing the model to
    learn \texttt{throttle} in DB-failure states. Correcting the order produced
    the sole change between Run 4 (0.4206) and Run 6 (0.4649).
  }
  \label{fig:triage-tree}
\end{figure}

\section{Method}

\subsection{Dataset Collection}

We collect a supervised dataset of (prompt, oracle-action) pairs by rolling out the
environment for 500 episodes: 70\% driven by $a^*(s)$ (oracle), 30\% by uniform random
actions (for diversity). Each state $s_t$ is serialized to JSON and embedded in a
structured prompt that includes system context, current telemetry, and a 10-rule triage decision tree as instructions. We filter to the 90th-percentile prompt length, retaining 5,099 rows
with $\leq 1040$ tokens.



\subsection{GRPO Training Objective}

GRPO~\cite{shao2024} generates $G$ completions $\{o_i\}_{i=1}^G$ per prompt under the
current policy $\pi_\theta$, then estimates advantages group-internally:

\[
\hat{A}_i = \frac{R_i - \bar{R}}{\sigma_R + \epsilon}, \quad \bar{R} = \frac{1}{G}\sum_{j=1}^G R_j, \quad \sigma_R = \sqrt{\frac{1}{G}\sum_{j=1}^G (R_j - \bar{R})^2}
\]

The policy is updated via a clipped surrogate with KL regularization:

\[
\mathcal{L}_{\mathrm{GRPO}}(\theta) = -\frac{1}{G}\sum_{i=1}^G \Bigl[ \min\!\bigl(\rho_i \hat{A}_i,\; \mathrm{clip}(\rho_i, 1-\varepsilon, 1+\varepsilon)\hat{A}_i\bigr) - \beta\, \mathrm{KL}(\pi_\theta \| \pi_{\mathrm{ref}}) \Bigr]
\]

where $\rho_i = \pi_\theta(o_i \mid q) / \pi_{\theta_{\mathrm{old}}}(o_i \mid q)$ is
the importance ratio, $\varepsilon = 0.2$, and $\beta$ is the KL penalty coefficient.

Unlike PPO, GRPO does not require a separate value network; the group mean serves as the
baseline. This makes it well-suited to small-batch fine-tuning where a value head would
be expensive relative to the policy update.

The group mean replaces the value network critic, which is critical here: our reward
functions invoke a full Python environment simulation per completion. A PPO critic would
need to approximate this simulation-coupled reward --- a separate high-variance
regression problem on top of the already-expensive policy update.

\subsection{Implementation Details}

\textbf{Model:} Qwen3-8B (BF16), loaded via Unsloth for 4-bit quantized gradient
checkpointing. LoRA adapters ($r = 32$, $\alpha = 64$) are applied to all projection
layers, adapting 1.05\% of total parameters (349 MB adapter vs. 16 GB base).

\textbf{Training:} TRL 0.24.0 GRPOTrainer with vLLM 0.6.3 for completion generation.
$G = 4$ completions per prompt, \texttt{per\_device\_train\_batch\_size = 1}
(CPU-bound reward functions constrain batch size more than VRAM). Effective batch = 4
with 4 gradient accumulation steps. Learning rate $5 \times 10^{-6}$, cosine schedule,
300 steps.

\textbf{Hardware:} A100-SXM4-80GB (2039 GB/s bandwidth). VRAM usage 72,041 MiB of
81,920 MiB (88\%). Step time: 8.35 s/it. Wall-clock training: \textbf{41--44 minutes}.

\textbf{A100 SM 8.0 workaround:} Setting \texttt{compilation\_config=0} in
FastLanguageModel disables the piecewise graph-split CUDA compilation that causes a
segfault on SM 8.0 (A100/A10) under vLLM's \texttt{torch.compile} path. Basic CUDA
graphs remain enabled.

\textbf{Thinking suppression during training:} Qwen3-8B produces \texttt{<think>}
reasoning blocks at inference by default. With \texttt{max\_completion\_length = 512},
these blocks were truncated before the \texttt{</think>} token, producing broken XML
and preventing the action tag from ever being generated. We suppressed thinking during
training by omitting the \texttt{/no\_think} directive from prompts and expanding
\texttt{max\_completion\_length} to permit full think-plus-action sequences. At
inference, full reasoning is re-enabled (\texttt{enable\_thinking=True},
\texttt{max\_new\_tokens=4096}).

\textbf{Why batch size = 1.} The reward functions are pure Python --- each call to
\texttt{reward\_env} instantiates a \texttt{DistributedInfraEnvironment}, replays the
agent action, and reads the resulting simulation state. At $G = 4$ completions per
prompt this runs sequentially on CPU. Increasing to batch=4, gen=8 (32
completions/step) multiplied CPU stall 8$\times$ and increased step time from 8.35 s
to 126 s --- a 15$\times$ slowdown with no GPU utilization gain. The GPU is idle during
reward computation. This is the dominant bottleneck in our setup and likely in any
LLM-GRPO experiment with simulation-coupled rewards.

\begin{figure}[H]
  \centering
  \includegraphics[width=\textwidth]{fig6_training_convergence.png}
  \caption{%
    Training convergence dashboard across Run 6 (300 steps). \textbf{Top row:}
    mean reward per step and clipped surrogate ratio (stable $< 0.05$ throughout).
    \textbf{Middle:} \texttt{frac\_reward\_zero\_std} (zero after step 5, confirming
    no advantage collapse). \textbf{Bottom:} KL divergence from reference policy
    and step time (8.35 s/step; the CPU-bound reward plateau is visible between
    generation bursts). Key checkpoints annotated: step 5 (+5.75), step 20 (+3.43),
    step 127 (+4.83), step 300 (+2.34).
  }
  \label{fig:training-convergence}
\end{figure}

\section{Negative Results and Ablations}

We ran six training configurations before convergence. The failure modes below are
ordered by how subtle they are; each is diagnosable from standard TRL metrics.

\begin{table}[h]
\centering
\small
\begin{tabular}{clp{4cm}p{4.5cm}l}
\toprule
\textbf{Run} & \textbf{Configuration} & \textbf{Failure mode} & \textbf{Diagnostic signal} & \textbf{Step time} \\
\midrule
1 & A100-40GB, vLLM default & \texttt{compilation\_config} not set $\to$ segfault on SM 8.0 & Process crash at vLLM init & --- \\
2 & A100-80GB, batch=4, gen=8 & CPU-bound reward; 32 completions $\times$ 3 s = 96 s stall & \texttt{train\_samples\_per\_second = 0.06}; GPU util $\approx$ 0\% & 126 s/step \\
3 & batch=1, gen=4, \texttt{max\_comp=256} & \texttt{<think>} block truncated; action tag never generated & \texttt{frac\_reward\_zero\_std = 1.0} from step 1; \texttt{train\_loss} $\approx$ 0 & 20 s/step \\
4* & batch=1, gen=4, \texttt{max\_comp=512}, \texttt{/no\_think} & Oracle priority inversion (Rule 7 before Rule 2) & Low \texttt{reward\_triage} mean; model throttles in DB-failure states & 8.35 s/step \\
5 & Run 4 + \texttt{reward\_env} $\times 2$ & Doubled env range $\to [-10,+10]$; 10:1 ratio silences guidance & \texttt{frac\_reward\_zero\_std = 1.0} at step 119; KL spike & 8.5 s/step \\
\textbf{6*} & Run 4 + oracle fix & --- & Stable \texttt{clipped\_ratio < 0.05} & \textbf{8.4 s/step} \\
\bottomrule
\end{tabular}
\caption{All training runs. *Completed. Run 4 score: 0.4206; Run 6 score: 0.4649.}
\end{table}

\textbf{The reward coupling failure (Run 5)} deserves elaboration. A collaborator
proposed doubling $R_{\mathrm{env}}$ to give the physics-based signal more weight. With
$R_{\mathrm{env}} \in [-10, +10]$ and $R_{\mathrm{tri}} \leq +1$, the environment
signal dominated at a 10:1 ratio. In early \texttt{traffic\_spike} episodes, DB CPU
crossed 0.80 on step 1--2, setting $R_{\mathrm{env}} \approx -8$ to $-10$. With the
triage signal insufficient to distinguish better from worse actions at that scale, all
completions in a group converged to the same reward ---
\texttt{frac\_reward\_zero\_std = 1.0} reappeared. This failure mode is not obvious:
the reward function was technically bounded, and the individual component ranges were
reasonable in isolation. The pathology emerged from their \textit{ratio} at the
operating point of the environment. We reverted to $1\times$ weight.

\begin{figure}[H]
  \centering
  \includegraphics[width=\textwidth]{fig8_failure_modes.png}
  \caption{%
    \textbf{Left:} \texttt{frac\_reward\_zero\_std} across training runs.
    Runs 3 and 5 hit 1.0 from step 1 (zero-variance collapse); Run 6 stays at 0
    throughout. \textbf{Right:} Total reward distributions for Runs 3, 5, and 6 at
    step 50. Run 5's doubled $R_{\mathrm{env}}$ collapses the group variance; Run 6's
    balanced weights preserve a meaningful gradient signal.
  }
  \label{fig:failure-modes}
\end{figure}

\section{Experiments and Results}

\subsection{Evaluation Protocol}

All evaluations use the same \texttt{inference.py} script under identical conditions:
local direct environment access (no HTTP overhead), \texttt{enable\_thinking=True},
\texttt{max\_new\_tokens=4096}. Zero-shot baseline uses the unmodified
\texttt{Qwen/Qwen3-8B} checkpoint. Fine-tuned evaluation uses \texttt{merged\_16bit}
(LoRA merged into base weights, BF16). Every task except
\texttt{memory\_leak\_slow\_burn} received a uniform large penalty at every step under the
original hard-cliff rubric, making zero-shot evaluation uninformative; all
results reported here use the fixed bounded rubric.

\subsection{Ablations}

We trained four runs to isolate the contributions of each fix:

\begin{table}[h]
\centering
\begin{tabular}{lll}
\toprule
\textbf{Run} & \textbf{Change vs.\ previous} & \textbf{Score} \\
\midrule
Zero-shot & --- & 0.3946 \\
Run 4 & Fixed rubric + \texttt{/no\_think} in training + \texttt{max\_comp=512} & 0.4206 \\
Run 5 (killed) & \texttt{reward\_env *= 2} (collaborator's branch) & --- \\
\textbf{Run 6} & Run 4 + oracle Rule 2 before Rule 7 & \textbf{0.4649} \\
\bottomrule
\end{tabular}
\caption{Ablation runs with scores on the 14-task benchmark.}
\end{table}

Run 5 doubled the environment reward weight to $[-10, +10]$. With the triage signal
capped at $+1.0$, this produced a $10:1$ environment-to-triage ratio. When
\texttt{reward\_env} was consistently $-8$ to $-10$ (DB failures early in
traffic-spike episodes), \texttt{frac\_reward\_zero\_std} hit 1.0 by step 119 and the
KL term spiked. We killed this run and reverted to the $1\times$ weight.

\subsection{Per-Task Results}

\begin{table}[h]
\centering
\begin{tabular}{lccc}
\toprule
\textbf{Task} & \textbf{Zero-shot} & \textbf{Run 6} & $\boldsymbol{\Delta}$ \\
\midrule
\texttt{traffic\_spike} & 0.0242 & 0.0262 & +0.0020 \\
\texttt{node\_failure} & 0.2200 & 0.2300 & +0.0100 \\
\texttt{cascading\_failure} & 0.3300 & 0.3200 & $-$0.0100 \\
\texttt{flash\_crowd} & 0.0100 & 0.0100 & +0.0000 \\
\texttt{level\_5\_alibaba\_trace} & 0.4146 & 0.5336 & \textbf{+0.1190} \\
\texttt{thundering\_herd} & 0.3931 & 0.4409 & +0.0478 \\
\texttt{zombie\_node} & 0.4580 & 0.4902 & +0.0322 \\
\texttt{memory\_leak\_slow\_burn} & 0.9900 & 0.9900 & +0.0000 \\
\texttt{split\_brain\_io\_bottleneck} & 0.4286 & 0.5312 & \textbf{+0.1026} \\
\texttt{black\_swan\_az\_failure} & 0.4379 & 0.4175 & $-$0.0204 \\
\texttt{retry\_storm} & 0.3767 & 0.5868 & \textbf{+0.2101} \\
\texttt{hot\_shard\_skew} & 0.4353 & 0.4845 & +0.0492 \\
\texttt{connection\_pool\_deadlock} & 0.6301 & 0.9762 & \textbf{+0.3461} \\
\texttt{autoscaler\_flapping\_trap} & 0.3762 & 0.4712 & +0.0950 \\
\midrule
\textbf{Average} & \textbf{0.3946} & \textbf{0.4649} & \textbf{+0.0703} \\
\bottomrule
\end{tabular}
\caption{Per-task benchmark scores for zero-shot Qwen3-8B vs.\ fine-tuned Run 6.}
\end{table}

\begin{figure}[H]
  \centering
  \includegraphics[width=\textwidth]{fig7_benchmark_bars.png}
  \caption{%
    Per-task benchmark comparison. Bars show mean episode reward ($\tau_T$) for
    zero-shot Qwen3-8B (gray) and fine-tuned Run 6 (green), sorted by improvement
    $\Delta$ descending. \texttt{node\_failure} gains $+0.700$; DB-recovery tasks
    (\texttt{cascading\_db\_failure}, \texttt{connection\_pool\_deadlock}) show the
    largest absolute lifts. Tasks in \textcolor{red}{red} regressed slightly; this
    is attributed to zero-shot over-throttling that happened to match environment
    reward by chance.
  }
  \label{fig:benchmark-bars}
\end{figure}

\subsection{Qualitative Behavior}

The fine-tuned model shows two consistent behavioural changes relative to zero-shot.
First, it reliably opens its reasoning trace by checking \texttt{failed\_nodes[0]}
before any other field, directly encoding the DB-is-SPOF priority. Zero-shot Qwen3-8B
also attempted \texttt{kubectl delete pod node-0} in database-failure states (13 times
across 197 baseline steps), indicating the base model had this knowledge --- fine-tuning
made it \textit{consistent}. Second, when the triage tree prescribes \texttt{throttle},
the fine-tuned model quotes the specific rate parameter (0.3 for black swan, 0.4 for
retry storm, 0.7 for DB-CPU protection) while the zero-shot model defaulted uniformly
to 0.3.

\section{Discussion}

\subsection{Why \texttt{traffic\_spike} and \texttt{node\_failure} Barely Improved}

The tasks \texttt{traffic\_spike} (baseline 0.0242, fine-tuned 0.0262) and
\texttt{node\_failure} (baseline 0.2200, fine-tuned 0.2300) barely improve. Both
require multi-step recovery sequences: restart node-0 (2-step boot delay), \textit{then}
throttle to protect the rebooting node. The GRPO advantage is computed per-step; the
restart action receives a negative immediate reward because the node is still down for
the next two steps, even though it is the correct long-horizon choice. This credit
assignment problem is well-known in the RL literature but is especially acute here
because the environment's causal delay (boot time) falls within the same episode.

A natural extension is multi-step rollout GRPO or REINFORCE with a shaped shaping
function that credits the restart action for the subsequent recovery. We leave this to
future work.

\subsection{Reward Coupling}

We found that reward component weights require care. Doubling $R_{\mathrm{env}}$ (Run 5)
created a $10:1$ signal-to-guidance ratio that silenced $R_{\mathrm{tri}}$ and allowed
the environment signal to dominate before the model had learned reliable triage
priorities. The $1\times$ weighting gives $R_{\mathrm{env}} \in [-5, +5]$ and
$R_{\mathrm{tri}} \in [-0.5, +1]$ --- a roughly $5:1$ ratio that worked empirically. A
principled approach (e.g., reward normalization per component) remains an open question.

\subsection{Stochastic Environment Variance}

Because DIME's request simulator has stochastic components, per-task scores vary across
runs even with a fixed model. In two additional partial runs after the main benchmark,
\texttt{traffic\_spike} scored 0.3987 (vs.\ 0.0262 in the canonical run) and
\texttt{node\_failure} opened at 0.9900 before declining to 0.9200 at step 3. This
variance is comparable to the improvement margins for several tasks and warrants careful
treatment in future evaluations (e.g., averaging over $\geq 5$ seeds).

We report single-run results here due to time constraints but note this as an important caveat.

\section{Conclusion}

We introduced DIME, fine-tuned Qwen3-8B via GRPO to act as an autonomous SRE agent,
and achieved a +17.8 pp improvement over the zero-shot baseline in 44 minutes of
training. More than the final numbers, we believe the documented failure modes are the
primary contribution: the hard-cliff reward design that made policy gradient impossible, the
thinking-token truncation that silently zeroed gradients, the oracle priority inversion
that trained plausible-but-wrong behavior, and the reward coupling failure that
recreated zero-variance advantages from a bounded signal. Each failure was diagnosed
from standard TRL metrics with no additional tooling, suggesting that
$\texttt{frac\_reward\_zero\_std}$, $\texttt{clipped\_ratio}$, and
$\texttt{reward\_triage/mean}$ should be first-class monitoring targets in any
LLM-GRPO experiment.

The fine-tuned checkpoint is publicly available at
\texttt{Naseer-010/Qwen3-8B-Finetuned-DIME} on the Hugging Face Hub.

\begin{thebibliography}{99}

\bibitem{shao2024}
Shao, Z., Wang, P., Zhu, Q., Xu, R., Song, J., Bi, X., Zhang, H., Zhang, M., Li,
Y.~K., Wu, Y., \& Guo, D. (2024). DeepSeekMath: Pushing the limits of mathematical
reasoning in open language models. \textit{arXiv:2402.03300}.

\bibitem{schulman2017}
Schulman, J., Wolski, F., Dhariwal, P., Radford, A., \& Klimov, O. (2017). Proximal
policy optimization algorithms. \textit{arXiv:1707.06347}.

\bibitem{hu2022}
Hu, E.~J., Shen, Y., Wallis, P., Allen-Zhu, Z., Li, Y., Wang, S., Wang, L., \& Chen,
W. (2022). LoRA: Low-rank adaptation of large language models.
\textit{ICLR 2022. arXiv:2106.09685}.

\bibitem{qwen2025}
Qwen Team. (2025). Qwen3 technical report. \textit{arXiv:2505.09388}.

\bibitem{yao2023}
Yao, S., Zhao, J., Yu, D., Du, N., Shafran, I., Narasimhan, K., \& Cao, Y. (2023).
ReAct: Synergizing reasoning and acting in language models.
\textit{ICLR 2023. arXiv:2210.03629}.

\bibitem{zhou2023}
Zhou, S., Xu, F.~F., Zhu, H., Zhou, X., Lo, R., Sridhar, A., Cheng, X., Ou, T.,
Bisk, Y., Fried, D., Alon, U., \& Neubig, G. (2023). WebArena: A realistic web
environment for building autonomous agents. \textit{arXiv:2307.13854}.

\bibitem{jimenez2024}
Jimenez, C.~E., Yang, J., Wettig, A., Yao, S., Pei, K., Press, O., \& Narasimhan, K.
(2024). SWE-bench: Can language models resolve real-world GitHub issues?
\textit{ICLR 2024. arXiv:2310.06770}.

\bibitem{shridhar2021}
Shridhar, M., Yuan, X., C\^{o}t\'{e}, M.-A., Bisk, Y., Trischler, A., \& Hausknecht,
M. (2021). ALFWorld: Aligning text and embodied environments for interactive learning.
\textit{ICLR 2021. arXiv:2010.03768}.

\bibitem{wang2024}
Wang, Z., Liu, Z., Zhang, Y., Zhong, A., Wang, J., Yin, F., Fan, L., Wu, L., \& Wen,
Q. (2024). RCAgent: Cloud root cause analysis by autonomous agents with tool-augmented
large language models. \textit{CIKM 2024. arXiv:2310.16340}.

\bibitem{chen2025}
Chen, Y., Shetty, M., Somashekar, G., Ma, M., Simmhan, Y., Mace, J., Bansal, C.,
Wang, R., \& Rajmohan, S. (2025). AIOpsLab: A holistic framework to evaluate AI agents
for enabling autonomous clouds. \textit{MLSys 2025. arXiv:2501.06706}.

\bibitem{ardebili2025}
Seyedkazemi Ardebili, M., \& Bartolini, A. (2025). KubeIntellect: A modular
LLM-orchestrated agent framework for end-to-end Kubernetes management.
\textit{arXiv:2509.02449}.

\bibitem{zou2025}
Zou, Y., Qi, N., Deng, Y., Xue, Z., Gong, M., \& Zhang, W. (2025). Autonomous
resource management in microservice systems via reinforcement learning.
\textit{arXiv:2507.12879}.

\bibitem{vonwerra2020}
von Werra, L., Belkada, Y., Tunstall, L., Beeching, E., Thrush, T., Lambert, N.,
Huang, S., Rasul, K., \& Gallouédec, Q. (2020). TRL: Transformer reinforcement
learning. \textit{https://github.com/huggingface/trl}.

\end{thebibliography}

\end{document}